
\documentclass[addpoints]{exam}
\usepackage[version=4]{mhchem}
\usepackage{siunitx}
\usepackage{color}
\usepackage{amsmath}

% \printanswers

\newcommand{\event}{Organic Chemistry}
\newcommand{\tournament}{}  % Replace with class name, date, or course code
\def\type{0}

\setlength\fillinlinelength{\textwidth}
\CorrectChoiceEmphasis{\color{red}\bfseries}
\SolutionEmphasis{\color{red}\bfseries}

\setlength\answerclearance{2pt}
\pagestyle{headandfoot}
\runningheadrule
\runningheader{\event}{\tournament}{\ifnum\type=1 Class Set Test \else Full Name:\hspace{1cm} \fi}
\runningfootrule
\runningfooter{}{Page \thepage\ of \numpages}{}

\begin{document}
\begin{center}
    {\huge Grade 12 Chemistry \\ \event
    \ifprintanswers
     \space Key
    \else
     \space \ifnum\type<2 Test \fi \ifnum\type=1 \textbf{(CLASS SET)} \fi \ifnum\type=2 Answer Sheet \fi
    \fi \\ \vspace{3mm}\tournament\vspace{1mm}}
\end{center}

\vspace{.25\textheight}

\begin{center}
    First Name: \ifprintanswers
    \underline{\hspace{3.5cm}}\textbf{KEY}\underline{\hspace{5.5cm}}\\\vspace{5mm}
    \else
        \underline{\hspace{10cm}}\\ \vspace{5mm}
    \fi
    Last Name: \underline{\hspace{10cm}}\\ \vspace{5mm}
\end{center}

\noindent\textbf{\underline{Directions:}}
\begin{itemize}
    \ifnum\type=1 \item \textbf{This test is a class set.} Please do not write any answers on this paper. \fi
    \item Show all work for full marks on short answer questions.
    \item Draw structural formulas clearly; use line structures where appropriate.
    \item The test is designed to be completed in 75 minutes.
\end{itemize}
\begin{center}
\textit{For grading use only}\\
\multirowgradetable{1}[pages]
\end{center}

\newpage
\begin{questions}

\section*{Part A --- Multiple Choice (15 marks)}
\vspace{5mm}

\question[1] What is the IUPAC name for the compound with the structure
\ce{CH3CH(CH3)CH2CH3}?
\begin{choices}
    \choice 2-methylpentane
    \correctchoice 2-methylbutane
    \choice 3-methylbutane
    \choice isopentane
\end{choices}

\question[1] Which functional group is present in a carboxylic acid?
\begin{choices}
    \choice \ce{-OH} only
    \choice \ce{C=O} only
    \correctchoice \ce{-COOH}
    \choice \ce{-NH2}
\end{choices}

\question[1] What is the IUPAC name for \ce{CH3COOC2H5}?
\begin{choices}
    \choice methyl propanoate
    \choice ethyl methanoate
    \correctchoice ethyl ethanoate
    \choice propyl methanoate
\end{choices}

\question[1] The addition of \ce{HBr} to an alkene is an example of which reaction type?
\begin{choices}
    \choice Substitution
    \choice Elimination
    \correctchoice Addition
    \choice Condensation
\end{choices}

\question[1] How many distinct structural isomers does \ce{C4H10} have?
\begin{choices}
    \choice 1
    \correctchoice 2
    \choice 3
    \choice 4
\end{choices}

\question[1] According to Markovnikov's rule, what is the major product when \ce{HBr} is added
to propene (\ce{CH3CH=CH2})?
\begin{choices}
    \choice 1-bromopropane
    \correctchoice 2-bromopropane
    \choice 1,2-dibromopropane
    \choice propan-1-ol
\end{choices}

\question[1] When a primary alcohol is oxidized under mild conditions, the initial product is a(n):
\begin{choices}
    \choice Ketone
    \choice Carboxylic acid
    \choice Ester
    \correctchoice Aldehyde
\end{choices}

\question[1] What is the organic product when ethanol undergoes dehydration (with \ce{H2SO4},
high temperature)?
\begin{choices}
    \choice Ethane
    \correctchoice Ethene
    \choice Ethanal
    \choice Diethyl ether
\end{choices}

\question[1] Which of the following is an addition polymer?
\begin{choices}
    \choice Nylon
    \choice Polyester
    \correctchoice Polyethylene
    \choice Protein
\end{choices}

\question[1] Saponification of a fat (triglyceride) with \ce{NaOH} produces:
\begin{choices}
    \choice An ester and water
    \choice An alcohol and a ketone
    \correctchoice Soap (fatty acid salts) and glycerol
    \choice An aldehyde and a carboxylic acid
\end{choices}

\question[1] What is the degree of unsaturation (index of hydrogen deficiency) of benzene,
\ce{C6H6}?
\begin{choices}
    \choice 2
    \choice 3
    \correctchoice 4
    \choice 5
\end{choices}

\question[1] The reaction \ce{CH4 + Cl2 ->[\text{UV}] CH3Cl + HCl} is classified as:
\begin{choices}
    \choice Electrophilic addition
    \choice Nucleophilic substitution
    \correctchoice Free-radical substitution
    \choice Elimination
\end{choices}

\question[1] Which condition is necessary for cis-trans (geometric) isomerism to exist in an
alkene?
\begin{choices}
    \choice A chiral carbon must be present.
    \correctchoice Each carbon of the double bond must have two different substituents.
    \choice The molecule must contain a ring.
    \choice The molecule must have at least four carbon atoms.
\end{choices}

\question[1] What are the two types of reactants required for an esterification reaction?
\begin{choices}
    \choice Aldehyde and alcohol
    \correctchoice Carboxylic acid and alcohol
    \choice Carboxylic acid and ketone
    \choice Alcohol and alkyl halide
\end{choices}

\question[1] What is the molecular formula of but-2-ene?
\begin{choices}
    \choice \ce{C4H6}
    \correctchoice \ce{C4H8}
    \choice \ce{C4H10}
    \choice \ce{C4H4}
\end{choices}


\section*{Part B --- Short Answer (33 marks)}

\question The molecular formula \ce{C5H12} has three structural isomers.
\begin{parts}
    \part[3] Draw and write the IUPAC name for all three structural isomers of \ce{C5H12}.
    \part[2] Which isomer has the lowest boiling point? Explain in terms of molecular structure
    and intermolecular forces.
    \part[3] The molecular formula \ce{C4H8} has several structural isomers including cyclic
    forms. Draw and name three distinct structural isomers of \ce{C4H8}, including at least
    one cyclic compound and one pair of cis/trans isomers.
\end{parts}
\begin{solution}[10cm]
\textbf{(a)}
\begin{enumerate}
    \item \textit{n}-pentane: \ce{CH3CH2CH2CH2CH3} (straight chain)
    \item 2-methylbutane: \ce{CH3CH(CH3)CH2CH3} (one branch)
    \item 2,2-dimethylpropane (neopentane): \ce{C(CH3)4} (two branches)
\end{enumerate}

\textbf{(b)} 2,2-dimethylpropane (neopentane) has the lowest boiling point (\SI{9.5}{\celsius},
vs.\ \SI{28}{\celsius} for 2-methylbutane and \SI{36}{\celsius} for \textit{n}-pentane). Its
compact, spherical shape gives a smaller surface area, resulting in weaker London dispersion forces
between molecules.

\textbf{(c)}
\begin{enumerate}
    \item but-1-ene: \ce{CH2=CHCH2CH3}
    \item \textit{cis}-but-2-ene: \ce{CH3CH=CHCH3} with same-side methyl groups
    \item \textit{trans}-but-2-ene: \ce{CH3CH=CHCH3} with opposite-side methyl groups
    \item Cyclobutane: four-carbon ring (also valid; any three of these four count)
\end{enumerate}
\end{solution}

\question The following questions concern reactions of alcohols.
\begin{parts}
    \part[4] Starting from propan-1-ol, write the structural formula of the organic product(s)
    formed in each of the following reactions. Name each product and identify the reaction type.
    \begin{itemize}
        \item[(i)] Oxidation with acidified \ce{K2Cr2O7} (excess oxidizing agent)
        \item[(ii)] Dehydration with concentrated \ce{H2SO4} at high temperature
        \item[(iii)] Reaction with ethanoic acid (\ce{CH3COOH}) in the presence of \ce{H2SO4}
        catalyst
    \end{itemize}
    \part[2] Distinguish between a primary, secondary, and tertiary alcohol. Give one example
    of each.
    \part[3] A secondary alcohol is treated with acidified \ce{K2Cr2O7}. Describe and explain
    what product forms, and contrast this with the oxidation of a primary alcohol.
\end{parts}
\begin{solution}[10cm]
\textbf{(a)}
\begin{itemize}
    \item[(i)] Excess oxidizing agent: propan-1-ol $\rightarrow$ propanal $\rightarrow$
    \textbf{propanoic acid} (\ce{CH3CH2COOH}). Reaction type: \textbf{oxidation}.
    \item[(ii)] \textbf{Propene} (\ce{CH3CH=CH2}). Reaction type: \textbf{elimination
    (dehydration)}.
    \item[(iii)] \textbf{Propyl ethanoate} (\ce{CH3COOCH2CH2CH3}) + \ce{H2O}. Reaction type:
    \textbf{esterification (condensation)}.
\end{itemize}

\textbf{(b)}
\begin{itemize}
    \item Primary: the carbon bearing \ce{-OH} is bonded to \textit{one} other carbon.
    Example: ethanol (\ce{CH3CH2OH}).
    \item Secondary: the carbon bearing \ce{-OH} is bonded to \textit{two} other carbons.
    Example: propan-2-ol (\ce{CH3CHOHCH3}).
    \item Tertiary: the carbon bearing \ce{-OH} is bonded to \textit{three} other carbons.
    Example: 2-methylpropan-2-ol (\ce{(CH3)3COH}).
\end{itemize}

\textbf{(c)} A secondary alcohol is oxidized to a \textbf{ketone} (e.g.\ propan-2-ol
$\rightarrow$ propanone). Oxidation stops at the ketone stage because the central carbonyl carbon
has no H atom to be removed in a further oxidation step. In contrast, a primary alcohol is
first oxidized to an \textbf{aldehyde}; with excess oxidizing agent, the aldehyde is further
oxidized to a \textbf{carboxylic acid}. A tertiary alcohol cannot be oxidized under these conditions
without breaking carbon--carbon bonds.
\end{solution}

\question The following questions concern polymers.
\begin{parts}
    \part[3] Distinguish between addition polymerization and condensation polymerization.
    Give one example of each type, identifying the monomer(s) and the polymer formed.
    \part[3] The following repeating unit belongs to a common polymer:
    \[
    \left[ \ce{-CH2-CH(Cl)-} \right]_n
    \]
    (i) Name this polymer and identify its monomer.\\
    (ii) Classify the polymerization as addition or condensation.\\
    (iii) Give one industrial use of this polymer.
    \part[2] Explain why condensation polymers such as nylon and polyester can be hydrolysed
    while addition polymers such as polyethylene cannot.
\end{parts}
\begin{solution}[9cm]
\textbf{(a)}
\begin{itemize}
    \item \textbf{Addition polymerization}: monomers containing a \ce{C=C} double bond add
    together; no small molecule is lost. Example: ethene (\ce{CH2=CH2}) $\rightarrow$
    polyethylene $[\ce{-CH2-CH2-}]_n$.
    \item \textbf{Condensation polymerization}: monomers join with loss of a small molecule
    (usually \ce{H2O}). Example: hexanedioic acid + 1,6-diaminohexane $\rightarrow$
    nylon-6,6 + \ce{H2O}.
\end{itemize}

\textbf{(b)}
\begin{itemize}
    \item[(i)] \textbf{Poly(chloroethene)}, commonly known as PVC (polyvinyl chloride).
    Monomer: chloroethene (\ce{CH2=CHCl}, vinyl chloride).
    \item[(ii)] \textbf{Addition} polymerization.
    \item[(iii)] Uses include pipes and plumbing, electrical cable insulation, window frames,
    flooring (any one acceptable).
\end{itemize}

\textbf{(c)} Condensation polymers contain hydrolysable linkages in the backbone
(\textbf{amide bonds} in nylon, \textbf{ester bonds} in polyester). Water can break these bonds
under acid or base catalysis, regenerating the original monomers. Addition polymers such as
polyethylene have only \ce{C-C} single bonds in the backbone, which are non-polar and highly
resistant to hydrolysis.
\end{solution}

\question An unknown organic compound X has the following properties:
\begin{itemize}
    \item Molecular formula \ce{C3H6O}
    \item Does not react with \ce{Na} metal to produce \ce{H2}
    \item Does not give a positive Tollens' test (silver mirror test)
    \item Burns completely in excess \ce{O2}
\end{itemize}
\begin{parts}
    \part[3] Using the above evidence, identify the functional group and name compound X.
    Justify each step of your reasoning.
    \part[3] Write a balanced equation for the complete combustion of compound X.
    \part[2] A structural isomer of X, compound Y (\ce{C3H6O}), gives a positive Tollens' test.
    Identify compound Y, draw its structural formula, and name the reaction type used in the
    Tollens' test.
\end{parts}
\begin{solution}[9cm]
\textbf{(a)} Reasoning:
\begin{itemize}
    \item Does not react with Na $\Rightarrow$ no \ce{-OH} group (rules out alcohol or
    carboxylic acid).
    \item Negative Tollens' test $\Rightarrow$ no aldehyde (\ce{-CHO}) group.
    \item Molecular formula \ce{C3H6O}: degree of unsaturation $= \frac{2(3)+2-6}{2} = 1$.
    One degree of unsaturation and a \ce{C=O} group with no OH and no aldehyde H $\Rightarrow$
    \textbf{ketone}.
\end{itemize}
Compound X is \textbf{propanone (acetone)}, \ce{CH3COCH3}.

\textbf{(b)}
\[
\ce{CH3COCH3 + 4O2 -> 3CO2 + 3H2O}
\]

\textbf{(c)} Compound Y gives a positive Tollens' test $\Rightarrow$ it contains an aldehyde
group. With formula \ce{C3H6O} and one degree of unsaturation: \textbf{propanal}
(\ce{CH3CH2CHO}).

The Tollens' test involves an \textbf{oxidation} reaction: the aldehyde is oxidized to a
carboxylate ion while \ce{Ag+} (in \ce{[Ag(NH3)2]+}) is reduced to silver metal (silver mirror).
\end{solution}

\end{questions}
\end{document}
